%=============================================================================
% ACTUAL MECHANISM: How the DFD Paper Produces alpha^{-1} = 137.03599985
%
% This document answers the question: what is the ACTUAL computation?
% Not what we wish it were. Not what a_4 should give. What it ACTUALLY does.
%
% Created: 2026-03-22
%=============================================================================
\documentclass[12pt,a4paper]{article}
\usepackage[margin=2.5cm]{geometry}
\usepackage{amsmath,amssymb}
\usepackage{booktabs,array,longtable}
\usepackage{hyperref}
\usepackage{tcolorbox}
\usepackage{xcolor}

\title{The Actual Mechanism Behind $\alpha^{-1} = 137.03599985$\\[6pt]
\large A Definitive Forensic Analysis of What the DFD Papers\\
Actually Compute, Step by Step}

\author{Forensic Audit --- All Source Documents Examined}
\date{March 22, 2026}

\begin{document}
\maketitle

\begin{abstract}
After exhaustive examination of three source documents ---
\texttt{section\_alpha\_locked.tex} (274 lines),
\texttt{appendix\_K.tex} (560+ lines), and the standalone PDF
``Ab Initio Derivation'' (27 pages, including Python code listings)
--- this document identifies the \textbf{actual mechanism} by which
DFD produces $\alpha^{-1} = 137.03599985$.

The answer is: \textbf{there are two distinct computations at two
different precision levels, connected by a formula that is never
explicitly written down}. The lattice route (Route~2) provides
$\sim 1\%$ verification. The sub-ppm value (Route~1) comes from
the spectral action framework, but the critical step --- the
computation of $\kappa_{U(1)}$ from the $a_4$ heat-kernel
coefficient --- is \textbf{never performed explicitly} in any
of the three source documents. This is the central gap.
\end{abstract}

\tableofcontents
\newpage

%=============================================================================
\section{Executive Summary: The Two-Route Architecture}
\label{sec:executive}
%=============================================================================

\begin{tcolorbox}[colback=red!5!white, colframe=red!60!black,
  title=Critical Finding]
The DFD corpus contains \textbf{two distinct derivation routes}
for $\alpha$:

\begin{enumerate}
\item \textbf{Route 1} (Section~8C / \texttt{section\_alpha\_locked.tex}):
  Claims sub-ppm precision ($\alpha^{-1} = 137.03599985$, residual
  $-0.006$~ppm). Lists eight algebraic/SM ingredients in a derivation
  chain table. \textbf{Never writes the master formula combining them.}

\item \textbf{Route 2} (Standalone PDF / \texttt{appendix\_K.tex}):
  Uses lattice Monte Carlo at $(\beta_{U(1)}, \beta_{SU(2)}) = (3.80, 22.80)$.
  Achieves $\sim 1\%$ precision. The formula is explicit and
  verifiable: Eq.~13 of the PDF.
\end{enumerate}

The sub-ppm claim comes \emph{exclusively} from Route~1.
Route~2 cannot deliver sub-ppm precision from lattice statistics.
\end{tcolorbox}

%=============================================================================
\section{What Route 2 Actually Computes (Verifiable)}
\label{sec:route2}
%=============================================================================

Route~2 is fully explicit. The standalone PDF provides two Python
code listings (pages 23--24) that implement the entire computation.

\subsection{Step 1: Compute $\beta_{U(1)}$ from CS Weights}

The weight function (Eq.~5 of PDF):
\begin{equation}
w(k) = \frac{2}{k+2}\sin^2\!\frac{\pi}{k+2}
\end{equation}
is the squared modulus of the SU(2) Chern--Simons partition function
$Z_{SU(2)_k}(S^3)$ on $S^3$ (Witten 1989).

The effective coupling (Eq.~15):
\begin{equation}
\beta_{U(1)} = \langle k_{\rm eff}\rangle
= \frac{\sum_{k=0}^{59}(k+2)\,w(k)}{\sum_{k=0}^{59}w(k)}
= 3.7969\ldots \approx 3.80
\end{equation}

This is a finite sum of trigonometric terms, computable to arbitrary
precision. The UV cutoff $k_{\max} = 60$ comes from the Bridge Lemma
(Eq.~14).

\subsection{Step 2: Derive Lattice Parameters}

\begin{align}
\beta_{SU(2)} &= \frac{n_2}{n_1} \times N_{\rm gen} \times \beta_{U(1)}
= 2 \times 3 \times 3.80 = 22.80 \tag{Eq.~17--18}
\end{align}

\subsection{Step 3: Run Lattice Monte Carlo}

At $(\beta_{U(1)}, \beta_{SU(2)}) = (3.80, 22.80)$, simulate compact
U(1) and SU(2) on an $L^4$ hypercubic lattice. Measure stiffnesses
$\kappa_{U(1)}$ and $\kappa_{SU(2)}$ via the background-field method
(Eq.~20):
\begin{equation}
\kappa = \frac{F''(0)}{V}, \qquad
F''(0) = \langle S''\rangle - \langle (S')^2\rangle + \langle S'\rangle^2
\end{equation}

\subsection{Step 4: Extract $\alpha$}

From measured stiffnesses (Eqs.~12--13):
\begin{align}
g_1^2 &= \frac{1}{\kappa_{U(1)}}, \qquad
g_2^2 = \frac{4}{\kappa_{SU(2)}} \\[6pt]
\alpha_W &= \frac{g_1^2 \cdot g_2^2}{g_1^2 + g_2^2} \cdot \frac{1}{4\pi}
\end{align}

\textbf{Python Listing 2} implements exactly this as
\texttt{alpha\_wilson(ku, ks)}.

\subsection{Step 4a: Simplification Using Stiffness Ratio}

With $\kappa_{SU(2)} = 2\kappa_{U(1)}$ (Frame Stiffness Theorem):
\begin{equation}
\boxed{\alpha^{-1} = 6\pi\,\kappa_{U(1)}}
\end{equation}

This is exact and proven. No approximation is involved.

\subsection{Route 2 Results}

\begin{center}
\begin{tabular}{lccc}
\toprule
$L$ & $\alpha_W$ (mean) & $\Delta\alpha/\alpha$ & Precision \\
\midrule
6 & 0.007297 & $-0.00\%$ & $\sim 1\%$ \\
8 & 0.007322 & $+0.34\%$ & $\sim 1\%$ \\
12 & 0.007291 & $-0.08\%$ & $\sim 0.3\%$ \\
\bottomrule
\end{tabular}
\end{center}

\textbf{Route 2 precision: $\sim 1\%$.} This is limited by lattice
statistics (finite $L$, finite sweeps, autocorrelations).

%=============================================================================
\section{What Route 1 Claims (Sub-ppm)}
\label{sec:route1}
%=============================================================================

Route~1 (\texttt{section\_alpha\_locked.tex}) claims
$\alpha^{-1} = 137.03599985$ with residual $-0.006$~ppm.

\subsection{The Derivation Chain Table}

Section~8C presents a table (Table~XXVIII) listing eight locked
ingredients. The table says ``All above combined'' gives
$\alpha^{-1} = 137.03599985$. But \textbf{no equation number,
no formula, and no computation is shown} for how these ingredients
combine to produce this number.

The ingredients are:
\begin{center}
\renewcommand{\arraystretch}{1.2}
\begin{tabular}{lll}
\toprule
Component & Value & Role \\
\midrule
$K_{\mathbb{C}P^2} = \mathcal{O}(-3)$ & $-3$ & Sets Spin$^c$ shift \\
$L_{\det} = K^{-1}$ & $\mathcal{O}(3)$ & Determinant line \\
$d = k+4$ & $64$ & Toeplitz dimension \\
$(d-1)/d$ & $63/64$ & Traceless projection \\
$N_{\rm species}$ & $7$ & SM content \\
$\text{Tr}(Y^2)$ & $10$ & Hypercharge sum \\
$g_F$ & $8$ & Spectral triple factor \\
$w = N_{\rm sp}/(g_F\text{Tr}(Y^2))$ & $7/80$ & Hypercharge weight \\
$\varepsilon_{\rm adj} = d^2/(d^2-1)$ & $4096/4095$ & BOOST (Branch A) \\
\bottomrule
\end{tabular}
\end{center}

\textbf{Notice:} $\beta_{U(1)} = 3.7969$ is listed in an updated version
of the table but was not in the original Section~8C table. The table
in Section~8C does \emph{not} list $\beta$ or $\kappa$ as ingredients.

\subsection{What Section 8C Actually Says}

The section makes the following claims:
\begin{enumerate}
\item The operator is the connection Laplacian on $\mathbb{C}P^2 \times S^3$
\item The regulator is a plateau cutoff (kills $a_6, a_8, \ldots$)
\item The truncation is Toeplitz at $d = 64$
\item The ``gauge-kinetic extraction from $a_4$ is unambiguous''
\item There is a forced binary fork between Branch A (BOOST) and Branch B (DROP)
\item Branch A gives 137.03599985; Branch B gives 137.03014445
\item Branch A matches experiment; Branch B is falsified at 43 ppm
\end{enumerate}

\textbf{What it does NOT do:}
\begin{itemize}
\item Write out the Gilkey formula for $a_4$
\item Compute any curvature invariant of $\mathbb{C}P^2$ or $S^3$
\item Show a product formula for $a_4$ on $M_1 \times M_2$
\item Cite Gilkey, Vassilevich, or Chamseddine--Connes
\item Display any intermediate value for $a_4$ or $\kappa$
\item Write a master formula combining the table ingredients
\end{itemize}

%=============================================================================
\section{The Central Gap: How Does $\beta = 3.80$ Become $\kappa = 7.27$?}
\label{sec:gap}
%=============================================================================

\begin{tcolorbox}[colback=orange!5!white, colframe=orange!60!black,
  title=The Missing Step]
The exact formula $\alpha^{-1} = 6\pi\kappa_{U(1)}$ requires
$\kappa_{U(1)} = 7.2700$ for the sub-ppm result.

The Chern--Simons computation gives $\beta_{U(1)} = 3.7969$.

The ratio $\kappa/\beta = 1.914$ must come from somewhere.

\textbf{The paper explicitly states (Eq.~11, page 7 of the PDF):}
``The relationship between [$\beta$ and $\kappa$] is non-perturbative
and lattice-size dependent, which is precisely why the simulation
is needed.''

Yet the sub-ppm value comes from Route~1, which does NOT use the
lattice simulation. The spectral action route must compute
$\kappa$ analytically --- but this computation is never shown.
\end{tcolorbox}

\subsection{What the Spectral Action Framework Should Provide}

In the Chamseddine--Connes spectral action approach:
\begin{equation}
S = \text{Tr}\,f(P/\Lambda^2) \sim \sum_{n \geq 0} f_n \cdot a_n(P)
\end{equation}

The $a_4$ coefficient contains the gauge-kinetic term:
\begin{equation}
a_4 \supset \frac{1}{(4\pi)^{n/2}} \int_M
\text{tr}\!\left(\frac{1}{12}\,\Omega_{ij}\Omega^{ij}\right)
\end{equation}
where $\Omega$ is the bundle curvature. This directly gives $\kappa$.

For the connection Laplacian on a line bundle $\mathcal{O}(1)$ over
$\mathbb{C}P^2 \times S^3$, evaluating $a_4$ requires:
\begin{itemize}
\item Curvature invariants of $\mathbb{C}P^2$ (never stated)
\item Curvature invariants of $S^3$ (never stated)
\item Metric normalization / volume of both spaces (never stated)
\item Product formula for $a_4$ on $M_1 \times M_2$ (never stated)
\item Toeplitz truncation effect on $a_4$ (never analyzed)
\end{itemize}

\textbf{None of these are provided in any of the three source documents.}

\subsection{Hypothesis Testing}

\paragraph{Hypothesis 1: The sub-ppm value comes from a specific lattice run.}
\textbf{Rejected.} The best lattice result (L12, $\alpha_W = 0.007291$,
deviation $-0.08\%$) has $\kappa_{U(1)} \approx 7.276$, not 7.270.
No lattice simulation at $L \leq 16$ can achieve sub-ppm precision.

\paragraph{Hypothesis 2: $\kappa/\beta \approx 1.91$ is computed from
compact U(1) perturbative expansion.}
\textbf{Not used.} The paper explicitly calls this relationship
``non-perturbative.'' Furthermore, the perturbative expansion of
$\langle P \rangle$ at $\beta = 3.80$ gives $\kappa/\beta \approx 0.93$
(from the plaquette mean field), not $1.91$.

\paragraph{Hypothesis 3: The sub-ppm value comes from a Chern--Simons
gauge coupling computation, not through $a_4$ at all.}
\textbf{Possible but unverifiable.} The paper mentions $a_4$
qualitatively but never computes it. The actual numerical path from
the CS weights to $\kappa = 7.270$ is not documented.

\paragraph{Hypothesis 4: CS/lattice duality maps $\beta$ directly to
$\kappa$.}
\textbf{No such duality is stated or used in the papers.}

%=============================================================================
\section{The Actual Answer: What Really Happens}
\label{sec:actual}
%=============================================================================

\begin{tcolorbox}[colback=blue!5!white, colframe=blue!60!black,
  title=Definitive Finding]
The DFD paper obtains $\alpha^{-1} = 137.03599985$ through the
following actual mechanism:

\begin{enumerate}
\item \textbf{The electroweak formula} $\alpha^{-1} = 6\pi\kappa_{U(1)}$
  is exact and proven (Frame Stiffness Theorem + Wilson normalization).

\item \textbf{The Chern--Simons weighted average}
  $\beta_{U(1)} = 3.7969$ is computable to arbitrary precision
  from the finite trigonometric sum with $k_{\max} = 60$.

\item \textbf{The lattice Monte Carlo} at $\beta = 3.80$ measures
  $\kappa_{U(1)} \approx 7.25 \pm 0.10$ (the PDF's Eq.~11 states
  this explicitly), confirming $\alpha \approx 1/137$ at $\sim 1\%$.

\item \textbf{The sub-ppm refinement} comes from Route~1
  (Section~8C), which claims that the spectral action on the
  Toeplitz-truncated $\mathbb{C}P^2 \times S^3$ computes
  $\kappa_{U(1)}$ analytically with all eight table ingredients.
  The BOOST factor $4096/4095$ is the final discrete correction
  that selects between Branch~A and Branch~B.

\item \textbf{The critical gap:} The formula that combines the eight
  ingredients to produce $\kappa_{U(1)} = 7.26999$ is never written
  down. Section~8C asserts the result but does not show the
  computation. The $a_4$ coefficient is invoked qualitatively
  (``the gauge-kinetic extraction from $a_4$ is unambiguous'')
  but never evaluated.
\end{enumerate}
\end{tcolorbox}

\subsection{Structural Diagnosis}

The derivation has three layers of decreasing explicitness:

\begin{center}
\renewcommand{\arraystretch}{1.3}
\begin{tabular}{lp{5cm}l}
\toprule
Layer & Content & Shown? \\
\midrule
$\alpha^{-1} = 6\pi\kappa_{U(1)}$ & Electroweak mixing formula &
  \textbf{Yes} (Eq.~13 of PDF) \\
$\kappa_{U(1)} \approx 7.25$ (lattice) & MC measurement at
  $\beta = 3.80$ & \textbf{Yes} (86 runs, Tables V--VIII) \\
$\kappa_{U(1)} = 7.26999\ldots$ (analytic) & Spectral action on
  Toeplitz $\mathbb{C}P^2 \times S^3$ &
  \textcolor{red}{\textbf{No}} (result stated, computation absent) \\
\bottomrule
\end{tabular}
\end{center}

The first two layers are fully explicit and verifiable.
The third layer --- which is the \emph{only} layer capable of
sub-ppm precision --- is asserted but not demonstrated.

\subsection{An Approximate Closed Form (85~ppm Residual)}

Reverse engineering suggests an approximate formula:
\begin{equation}
\alpha^{-1} \approx \frac{35\pi}{3}\;\beta_{U(1)}
\;\frac{d-1}{d}\;\frac{d^2}{d^2-1}
\end{equation}
which gives $\alpha^{-1} \approx 137.024$ (85~ppm from experiment).
The prefactor $35\pi/3 = 4\pi \times (5/2) \times (7/6)$ can be
decomposed into electroweak mixing and spectral triple factors.
But this is a \emph{reverse-engineered approximation}, not a
derived result.

%=============================================================================
\section{What Would Be Needed for a Complete Derivation}
\label{sec:needed}
%=============================================================================

To close the gap in Section~\ref{sec:gap}, one would need to:

\begin{enumerate}
\item \textbf{Compute $a_4$ explicitly} for the connection Laplacian
  on a line bundle over $\mathbb{C}P^2 \times S^3$, using the
  standard Gilkey formula with specified metric normalizations.

\item \textbf{Apply the Toeplitz truncation} and show how it
  modifies the $a_4$ coefficient (if at all --- previous analysis
  showed Toeplitz truncation does not modify polynomial SDW
  coefficients at leading order).

\item \textbf{Incorporate the SM content} ($N_{\rm species} = 7$,
  $\text{Tr}(Y^2) = 10$, $g_F = 8$) into the $a_4$ trace and
  show that the result gives
  $\kappa_{U(1)} = 137.03599985/(6\pi) = 7.26999$.

\item \textbf{Or:} Provide an alternative analytic formula
  (not going through $a_4$) that produces $\kappa = 7.27$ from
  $\beta = 3.80$ and the discrete factors.
\end{enumerate}

Previous attempts to compute $a_4$ directly have produced:
\begin{itemize}
\item Scalar Laplacian $a_4$: gives $\alpha^{-1} \approx 20$
  (factor 6.79 too small)
\item Dirac operator $a_4$: even smaller than scalar
\item Various prefactors ($35/18$, etc.): reverse-engineered,
  not derived from Gilkey
\end{itemize}

%=============================================================================
\section{Conclusions}
\label{sec:conclusions}
%=============================================================================

\begin{enumerate}
\item The sub-ppm value $\alpha^{-1} = 137.03599985$ is stated but
  not derived in the extant source documents. The derivation chain
  table lists ingredients; the formula combining them is absent.

\item The $\sim 1\%$-level result ($\alpha \approx 1/137$ from
  lattice MC) \emph{is} fully explicit and reproducible.
  The code is published (Listings 1--2 of the PDF; Zenodo repository).

\item The relationship $\beta \to \kappa$ (bare coupling to
  renormalized stiffness) is explicitly called ``non-perturbative''
  by the paper itself (Eq.~11). Yet Route~1 claims analytic
  sub-ppm precision for this very quantity without showing how.

\item The BOOST factor $4096/4095$ is well-defined and enters
  cleanly. It is the \emph{only} component of the sub-ppm claim
  that is fully derived and explicit.

\item \textbf{The actual mechanism is:} $\alpha^{-1} = 6\pi\kappa_{U(1)}$
  (exact), with $\kappa_{U(1)}$ measured by lattice MC at $\sim 1\%$
  precision. The sub-ppm claim requires an analytic computation of
  $\kappa$ that has not been shown.

\item None of the four hypotheses tested (lattice fine-tuning,
  perturbative $\kappa/\beta$, direct CS gauge coupling, CS/lattice
  duality) successfully explains the sub-ppm value from first
  principles.
\end{enumerate}

\end{document}
